\chapter{Diferencia Temporal: SARSA y Q-Learning}
\label{cap:td}

\section{La Gran Idea: El Aprendizaje por Bootstrapping}
El Aprendizaje por Diferencia Temporal (\textit{Temporal-Difference Learning}, TD) representa la idea más distintiva y profunda del Aprendizaje por Refuerzo. Descrita por Richard Sutton en 1988, la diferencia temporal fusiona magistralmente las mejores virtudes de los dos mundos estudiados en los capítulos anteriores:
\begin{itemize}
  \item Al igual que los \textbf{Métodos Monte Carlo}, TD aprende directamente a partir de la experiencia bruta en línea, sin necesidad de conocer el modelo de transición del entorno.
  \item Al igual que la \textbf{Programación Dinámica}, TD actualiza sus estimaciones basándose en parte en otras estimaciones previas, sin esperar al desenlace final del episodio. Esta propiedad se conoce formalmente como \textbf{bootstrapping}.
\end{itemize}

\begin{optibot}
OptiBot ya no espera a que termine su jornada laboral de 8 horas para calcular los retornos y corregir sus decisiones. Ahora, tras cada metro recorrido y cada segundo transcurrido, observa inmediatamente la recompensa obtenida $R_{t+1}$ y el nuevo estado $S_{t+1}$. Si el nuevo estado parece más ventajoso de lo previsto, incrementa al instante el valor del estado anterior; si detecta peligro o demora, reduce su estimación sin dilación.
\end{optibot}

\section{Predicción con TD(0)}
Mientras que Monte Carlo espera al retorno real completo $G_t$ para realizar la actualización:
\[
  V(S_t) \leftarrow V(S_t) + \alpha \left[ G_t - V(S_t) \right]
\]
el método más simple de Diferencia Temporal, denominado \textbf{TD(0)}, sustituye el retorno futuro desconocido por una aproximación en un solo paso:
\begin{equation}
  V(S_t) \leftarrow V(S_t) + \alpha \left[ \underbrace{R_{t+1} + \gamma V(S_{t+1})}_{\text{Objetivo TD}} - V(S_t) \right]
\end{equation}

La cantidad entre corchetes se denomina el \textbf{Error de Diferencia Temporal}, comúnmente abreviado como \textbf{Error TD} ($\delta_t$):
\begin{equation}
  \delta_t \doteq R_{t+1} + \gamma V(S_{t+1}) - V(S_t)
\end{equation}
El error TD cuantifica la discrepancia entre la estimación que el agente tenía en el paso $t$ y la evidencia mejorada obtenida tras avanzar un paso hacia $t+1$.

\section{Algoritmos de Control: SARSA vs. Q-Learning}
Para controlar el movimiento de OptiBot sin un modelo físico del almacén, debemos estimar las funciones de valor de acción $Q(s, a)$. Existen dos enfoques fundamentales según la relación entre la política que genera los datos y la política que se optimiza:

\subsection{1. SARSA: Control On-Policy}
El nombre SARSA proviene de la tupla de 5 elementos que intervienen en cada actualización: $(S_t, A_t, R_{t+1}, S_{t+1}, A_{t+1})$. La regla de actualización es:
\begin{equation}
  Q(S_t, A_t) \leftarrow Q(S_t, A_t) + \alpha \left[ R_{t+1} + \gamma Q(S_{t+1}, A_{t+1}) - Q(S_t, A_t) \right]
\end{equation}
En SARSA, la siguiente acción $A_{t+1}$ se selecciona utilizando la \textbf{misma política de comportamiento} $\varepsilon$-greedy con la que el agente está actuando. Por ello, SARSA aprende el valor de la política que realmente ejecuta, internalizando el riesgo de que el robot pueda cometer un error exploratorio aleatorio. En entornos con precipicios o riesgos graves, SARSA aprende rutas más \textbf{prudentes y conservadoras}.

\subsection{2. Q-Learning: Control Off-Policy}
Propuesto por Christopher Watkins en 1989, \textbf{Q-Learning} es uno de los algoritmos más influyentes de la inteligencia artificial. Su regla de actualización es:
\begin{equation}
  Q(S_t, A_t) \leftarrow Q(S_t, A_t) + \alpha \left[ R_{t+1} + \gamma \max_{a' \in \mathcal{A}} Q(S_{t+1}, a') - Q(S_t, A_t) \right]
\end{equation}

\begin{ideaclave}
Nótese la diferencia fundamental entre ambos algoritmos:
\begin{itemize}
  \item En SARSA, el objetivo utiliza $Q(S_{t+1}, A_{t+1})$, donde $A_{t+1}$ es la acción realmente seleccionada (que puede ser un movimiento exploratorio aleatorio y subóptimo).
  \item En Q-Learning, el objetivo utiliza $\max_{a'} Q(S_{t+1}, a')$, asumiendo que en el siguiente estado el agente actuará de forma estrictamente codiciosa (\textit{greedy}).
\end{itemize}
Esto permite a Q-Learning aprender la \textbf{política óptima directa $Q^*$} de forma \textit{off-policy}, con independencia de los desvíos exploratorios que el robot cometa durante el entrenamiento.
\end{ideaclave}

\begin{calculomanual}
Supongamos que OptiBot se encuentra en el Pasillo A ($S=0$) y ejecuta la acción Avanzar ($A=1$). Recibe una recompensa inmediata $R = -1.0$ y llega al Pasillo B ($S'=1$).
\begin{itemize}
  \item Valor actual en la tabla: $Q(0, 1) = 3.50$.
  \item En el Pasillo B, las estimaciones para sus 3 acciones son:
  \[
    Q(1, \text{GirarIzquierda}) = 6.00, \quad Q(1, \text{Avanzar}) = 8.00, \quad Q(1, \text{GirarDerecha}) = 4.00
  \]
  \item Parámetros de aprendizaje: tasa de paso $\alpha = 0.20$, factor de descuento $\gamma = 0.90$.
\end{itemize}
Desarrollamos el paso de actualización de Q-Learning con lápiz y papel:
\begin{enumerate}
  \item Calculamos el mejor valor de acción en el estado siguiente:
  \[
    \max_{a'} Q(S', a') = \max(6.00, \, 8.00, \, 4.00) = \mathbf{8.00}
  \]
  \item Construimos el objetivo de Bellman TD:
  \[
    \text{Objetivo} = R + \gamma \max_{a'} Q(S', a') = -1.0 + (0.90)(8.00) = -1.0 + 7.20 = \mathbf{6.20}
  \]
  \item Calculamos el error TD de sorpresa:
  \[
    \delta = \text{Objetivo} - Q(S, A) = 6.20 - 3.50 = \mathbf{+2.70}
  \]
  \item Actualizamos la celda correspondiente en la tabla Q:
  \[
    Q_{\text{nuevo}}(0, 1) = Q(0, 1) + \alpha \cdot \delta = 3.50 + (0.20)(+2.70) = 3.50 + 0.54 = \mathbf{4.04}
  \]
\end{enumerate}
La estimación del robot sube de $3.50$ a $4.04$, reflejando que el Pasillo B ofrece opciones de navegación altamente ventajosas.
\end{calculomanual}

\section{Implementación en Python: Q-Learning Tabular}
\begin{lstlisting}[style=rlcode]
import numpy as np

class QLearningAgent:
    def __init__(self, n_states, n_actions, alpha=0.1, gamma=0.95, epsilon=0.1):
        self.n_states = n_states
        self.n_actions = n_actions
        self.alpha = alpha
        self.gamma = gamma
        self.epsilon = epsilon
        self.Q = np.zeros((n_states, n_actions))

    def select_action(self, state):
        if np.random.rand() < self.epsilon:
            return np.random.randint(self.n_actions)
        return np.argmax(self.Q[state])

    def update(self, s, a, r, s_next, done):
        best_next_q = 0.0 if done else np.max(self.Q[s_next])
        td_target = r + self.gamma * best_next_q
        td_error = td_target - self.Q[s, a]
        self.Q[s, a] += self.alpha * td_error

# Bucle principal de interaccion con el entorno de OptiBot
# env = WarehouseEnv()
# for ep in range(1000):
#     s, _ = env.reset()
#     done = False
#     while not done:
#         a = agent.select_action(s)
#         s_next, r, done, _, _ = env.step(a)
#         agent.update(s, a, r, s_next, done)
#         s = s_next
\end{lstlisting}

\begin{notabox}
\textbf{La limitación que motiva el siguiente capítulo:} Q-Learning tabular es extraordinariamente robusto para problemas discretos pequeños. Sin embargo, su talón de Aquiles es la \textbf{maldición de la dimensionalidad}. Si el almacén de OptiBot mide $100 \times 100$ metros y el robot añade sensores continuos de distancia LiDAR, orientación continua en radianes $\theta \in [0, 2\pi)$ y lecturas de nivel de batería en tanto por ciento ($0.0\%$ a $100.0\%$), el número de estados posibles es literalmente \textbf{infinito}. Una tabla matricial no cabe en ninguna memoria RAM y jamás visitaría cada celda. Necesitamos pasar del almacenamiento tabular a la \textbf{Aproximación de Funciones con Redes Neuronales Profundas (DQN)}.
\end{notabox}
